\chapter{Conclusion générale}

\section{Bilan de l'étude}

Ce projet informatique avait pour objectif de développer un système d'estimation de la valeur
des biens immobiliers résidentiels au Sénégal, en explorant l'apport de l'optimisation des
hyperparamètres d'entraînement et en rendant les résultats accessibles via une interface
graphique interactive. Les trois objectifs fixés en introduction ont été atteints :

\begin{itemize}
  \item un jeu de données de 603 annonces exploitables a été constitué par extraction
  automatisée depuis Expat-Dakar, en l'absence de toute source structurée préexistante pour le
  marché sénégalais ;
  \item quatre modèles ont été entraînés et comparés selon un protocole homogène, avec
  optimisation systématique des hyperparamètres par recherche aléatoire et validation croisée
  pour les trois modèles les plus complexes ; le modèle Random Forest optimisé obtient les
  meilleures performances (MAE de 73,3 millions FCFA, MAPE de 34,0~\%, $R^2$ de 0,751),
  devançant XGBoost et le réseau de neurones MLP ;
  \item une application Streamlit interactive, déployée publiquement à l'adresse
  \url{https://estimation-valeurs-immobilieres.streamlit.app/} (code source disponible sur
  \url{https://github.com/Laminito/MemoireM1DIT.git}), permet d'estimer le prix d'un bien, de
  comparer les prédictions des quatre modèles, et d'explorer visuellement le jeu de données et
  les performances des modèles.
\end{itemize}

Le projet confirme, dans le contexte spécifique et peu documenté du marché immobilier
sénégalais, des tendances déjà observées dans la littérature internationale : la supériorité
des méthodes ensemblistes à base d'arbres sur la régression linéaire classique, et la nécessité
d'un volume de données conséquent pour que les réseaux de neurones expriment un avantage sur
ces méthodes.

\section{Perspectives}

Plusieurs axes d'amélioration se dégagent de ce travail :

\begin{enumerate}
  \item \textbf{Enrichissement des données.} L'intégration de sources complémentaires (autres
  plateformes d'annonces, données notariales si elles venaient à être ouvertes) permettrait
  d'augmenter le volume et la diversité géographique du jeu de données, actuellement très
  concentré sur Dakar.
  \item \textbf{Enrichissement des variables.} L'ajout de caractéristiques non disponibles dans
  les annonces actuelles (âge du bien, statut foncier, proximité d'axes ou d'équipements),
  à obtenir par géocodage des adresses, permettrait de rapprocher les
  performances du modèle de celles obtenues sur des jeux de données occidentaux plus riches en
  variables.
  \item \textbf{Suivi temporel des prix.} La collecte périodique des mêmes annonces
  permettrait de suivre l'évolution des prix dans le temps et d'enrichir le modèle d'une
  dimension temporelle, actuellement absente.
  \item \textbf{Modèles d'ensemble combinés.} Une combinaison (\emph{stacking}) des prédictions
  de Random Forest, XGBoost et du MLP pourrait améliorer encore la robustesse des estimations.
  \item \textbf{Déploiement public.} La mise en ligne de l'application sur Streamlit Community
  Cloud rendrait l'outil accessible au grand public, à l'image du projet de référence ayant
  inspiré la structure de ce document.
\end{enumerate}

En définitive, ce projet constitue une première contribution méthodologique à l'estimation
automatisée des prix immobiliers au Sénégal, et pose les bases d'un outil dont la précision
pourrait être significativement améliorée par un enrichissement futur des données mobilisées.
